\documentclass[10pt,twocolumn]{article}

% ── Packages ──────────────────────────────────────────────────────────────────
\usepackage[margin=1in]{geometry}
\usepackage{times}
\usepackage{graphicx}
\usepackage{amsmath,amssymb}
\usepackage{booktabs}
\usepackage{hyperref}
\usepackage{xcolor}
\usepackage{subcaption}
\usepackage{microtype}
\usepackage{cite}
\usepackage{algorithm}
\usepackage{algpseudocode}

% ── Metadata ──────────────────────────────────────────────────────────────────
\title{Shadow-Conditioned Tree Crown Detection for Carbon Estimation\\
       in Australian Arid Rangeland from Sub-Decimetre UAV Imagery}

\author{
  Tom Pitts$^{1}$ \\
  $^{1}$Department of Computer Science, University of Oxford \\
  \texttt{thomas.pitts@cs.ox.ac.uk}
}

\date{}

% ──────────────────────────────────────────────────────────────────────────────
\begin{document}
\maketitle

% ── Abstract ──────────────────────────────────────────────────────────────────
\begin{abstract}
Accurate estimation of above-ground biomass in Australian arid rangelands
requires reliable detection of sparsely distributed, isolated trees at scale.
We address this with a three-stage aerial Tree Crown Detection (TCD) pipeline
operating on proprietary 0.1\,m\,px$^{-1}$ UAV imagery: (1) shadow direction
prediction via a ResNet34 trained to regress sun azimuth directly from RGB,
(2) shadow-conditioned TCD using a Faster R-CNN with ResNet50-FPN backbone,
and (3) crown area estimation via Segment Anything Model (SAM) post-processing.
The central challenge in arid rangeland is visual ambiguity: sparse tree canopy
is nearly indistinguishable from low scrub and pale sandy substrate in RGB alone.
We exploit cast shadows as a reliable indicator of vertical extent --- a cue that
unambiguously separates trees from ground-level vegetation. We investigate
multiple mechanisms for incorporating shadow information into the detector,
finding that appending a predicted shadow probability map as a fourth input
channel consistently improves mean Average Precision (mAP@0.4) by
\textbf{+4.8\%} over the RGB-only baseline (0.650 vs.\ 0.620). Through an
ablation study isolating the contribution of the shadow \emph{direction} vector,
we find that directional and isotropic shadow maps produce equivalent detection
performance, suggesting the detector exploits shadow as an occlusion-awareness
signal rather than a geometric prior. These results motivate future work on
architectures that can explicitly reason about shadow geometry.
\end{abstract}

% ── 1. Introduction ───────────────────────────────────────────────────────────
\section{Introduction}

Australia's rangelands cover approximately 70\% of the continent and store
significant quantities of above-ground carbon in isolated trees and sparse
woodland~\cite{ma2021global}. Quantifying this carbon stock at national scale
requires accurate tree-level detection and crown delineation over vast areas ---
a task for which UAV remote sensing at sub-decimetre resolution is increasingly
practical but computationally demanding.

State-of-the-art TCD models trained on dense temperate forest imagery (e.g.,
NEON~\cite{weinstein2020neon}, NeonTreeEvaluation~\cite{weinstein2021benchmark})
generalise poorly to arid Australian rangeland. The domain gap is severe:
canopy density is an order of magnitude lower, the substrate is pale sandy soil
rather than dark earth or litter, and tree crowns at height 3--15\,m cast
long, clearly defined shadows on bare ground. Existing models trained without
this cue confuse open scrub, soil patches, and rock with tree crowns.

We make the following contributions:
\begin{enumerate}
  \item A \textbf{shadow direction prediction model}: a ResNet34 fine-tuned to
        regress sun azimuth from a single aerial RGB tile, enabling shadow map
        generation without GPS/IMU metadata.
  \item A \textbf{shadow-conditioned Faster R-CNN} that appends a classical
        shadow probability map as a fourth input channel to the backbone,
        improving recall on shadow-obscured crowns.
  \item A \textbf{three-stage detection pipeline} integrating shadow direction,
        TCD, and SAM-based segmentation for per-tree crown area estimation.
  \item An \textbf{ablation study} dissecting whether detection gains are
        attributable to the directional shadow vector or to a direction-free
        darkness cue, with implications for future geometric architectures.
\end{enumerate}

% ── 2. Related Work ───────────────────────────────────────────────────────────
\section{Related Work}

\paragraph{Tree Crown Detection.}
Deep learning TCD has been advanced primarily through object detection frameworks
applied to aerial and satellite imagery~\cite{weinstein2020neon,li2023single}.
DeepForest~\cite{weinstein2020neon} provides a Faster R-CNN baseline pre-trained
on NEON LiDAR-derived crowns and has become a standard starting point for
fine-tuning. Transformer-based detectors~\cite{li2023single} and instance
segmentation models~\cite{kirillov2023sam} have recently extended performance
on dense canopy but are not specifically designed for sparse arid domains.

\paragraph{Shadow in Remote Sensing.}
Shadows in aerial imagery are typically treated as nuisances to be suppressed
during pre-processing~\cite{adeline2013shadow}. Shadow removal~\cite{liu2020shadow}
and shadow detection~\cite{wang2021shadow} have received significant attention in
the street-view and satellite domains, but shadow as a \emph{positive detection
cue} has not been systematically exploited in TCD. Physiologically, a cast shadow
indicates vertical extent above ground and therefore reliably distinguishes trees
from low scrub --- this is the key insight we develop.

\paragraph{Sun Azimuth Estimation.}
Sun direction can be derived from GPS coordinates, acquisition time, and
ephemeris~\cite{reda2004solar} when metadata is available. For proprietary UAV
datasets where flight metadata is restricted, we instead learn to predict azimuth
directly from image statistics (shadow orientation, highlight gradients), following
the approach of~\cite{lalonde2010estimating} for outdoor scene understanding.

\paragraph{Rangeland Carbon.}
Remote sensing-based biomass estimation in Australian rangelands has relied on
coarse satellite imagery ($\ge$10\,m) and allometric relationships~\cite{ma2021global}.
Sub-decimetre UAV surveys enable individual tree-level detection, opening the
possibility of trunk-diameter inference from crown area and direct integration
with allometric equations for carbon estimation.

% ── 3. Method ─────────────────────────────────────────────────────────────────
\section{Method}

The pipeline comprises three sequential stages, each consuming the outputs of
the previous. An overview is shown in Figure~\ref{fig:pipeline}.

\begin{figure}[t]
  \centering
  % TODO: replace with actual pipeline figure
  \fbox{\rule{0pt}{4cm}\rule{\linewidth}{0pt}}
  \caption{Three-stage pipeline. Stage 1 predicts shadow direction from RGB.
           Stage 2 runs shadow-conditioned TCD. Stage 3 refines crown polygons
           with SAM. Red boxes are detected crowns; blue overlays are SAM masks.}
  \label{fig:pipeline}
\end{figure}

\subsection{Stage 1: Shadow Direction Prediction}

Given a UAV tile $\mathbf{I} \in \mathbb{R}^{H \times W \times 3}$, we predict
the sun azimuth $\phi \in [0°, 360°)$ as a unit vector
$\mathbf{s} = (\sin\phi, \cos\phi)$, avoiding the angular discontinuity at $0°$.

A ResNet34 backbone is fine-tuned with a two-output regression head minimising
the mean angular error:
\begin{equation}
  \mathcal{L}_\phi = \arccos\!\left(
    \hat{\mathbf{s}} \cdot \mathbf{s}^* / \|\hat{\mathbf{s}}\| \|\mathbf{s}^*\|
  \right),
\end{equation}
where $\mathbf{s}^*$ is the ground-truth direction derived from GPS and
ephemeris~\cite{reda2004solar}. At inference the model produces a per-tile
direction without requiring any metadata.

\subsection{Stage 2: Shadow-Conditioned Tree Crown Detection}
\label{sec:stage2}

\paragraph{Shadow Probability Map.}
Given $\mathbf{I}$ and azimuth $\phi$, we compute a shadow probability map
$\mathbf{M} \in [0,1]^{H \times W}$ using a classical directional gradient
algorithm. For each pixel $p$, we measure how much brighter nearby pixels are
in the sun direction (sun-side maximum) relative to the anti-sun direction:
\begin{equation}
  \mathrm{dg}(p) = \left[B_{\mathrm{sun}}(p) - B(p)\right]_+ -
                    0.4\left[B_{\mathrm{anti}}(p) - B(p)\right]_+,
\end{equation}
where $B_{\mathrm{sun}}(p) = \max_{d \in \mathcal{D}} B(p + d\hat{\mathbf{s}})$
accumulates the maximum brightness seen at offsets $d \in \mathcal{D}$
(short range: 8--35\,px; long range: 35--150\,px) in the sun direction, and
$[\cdot]_+ = \max(\cdot, 0)$.
The anti-sun term penalises globally dark pixels where no directional asymmetry
exists.

A luma gate $\lambda(p) = \operatorname{clip}\!\left(\frac{L_{\max} - B(p)}{20}, 0, 1\right)$
suppresses pixels brighter than domain-specific threshold $L_{\max}$
(BRU: 71; WON: 150, reflecting brighter sandy substrate).
A lateral Sobel penalty suppresses road edges parallel to the shadow direction.
The combined activation is passed through a sigmoid with fixed centre 0.35 and
speckle-removal ($\ge 150$\,px$^2$) to produce $\mathbf{M}$.

\paragraph{Shadow-Conditioned Detector.}
We fine-tune DeepForest~\cite{weinstein2020neon}, a Faster R-CNN with
ResNet50-FPN backbone pre-trained on NEON aerial data.
In our shadow-conditioned variant, $\mathbf{M}$ is appended to the RGB image
as a fourth channel before the first convolutional layer, which is widened
from $3 \to 4$ input filters with zero initialisation of the new weights.
This preserves the pre-trained RGB representation at initialisation while
allowing the backbone to learn to integrate shadow evidence from the first
layer through a natural hierarchical feature representation.

\paragraph{WON Annotation Normalisation.}
Ground-truth bounding boxes in WON imagery (site introduced in
Section~\ref{sec:datasets}) were annotated to include the cast shadow in
addition to the canopy, producing boxes $\approx 2\times$ larger in area
than equivalent BRU annotations. We correct for this at training time by
shrinking WON boxes to 50\% of their original area (scaling each side by
$\sqrt{0.5} \approx 0.707$) around the box centre without modifying the
source CSVs.

\subsection{Stage 3: Crown Segmentation via SAM}

Detected bounding boxes from Stage 2 are passed as prompts to the Segment
Anything Model (SAM)~\cite{kirillov2023sam} to produce per-crown segmentation
masks. Crown pixel area is converted to physical area at the known ground
sampling distance (0.1\,m\,px$^{-1}$) and used as input to published allometric
equations for estimating above-ground biomass and carbon stock.

% ── 4. Experimental Setup ─────────────────────────────────────────────────────
\section{Experimental Setup}

\subsection{Datasets}
\label{sec:datasets}

We evaluate on two proprietary UAV datasets acquired at 0.1\,m\,px$^{-1}$
over Australian rangeland:

\textbf{WON} (\emph{Wongan Hills, Western Australia}): Sandy pale substrate,
sparse eucalyptus woodland with tree density 5--30 trees/ha. Trees are
well-separated, casting long shadows on bare ground. Shadow luma values
(80--140) are substantially higher than darker-substrate sites, requiring
a relaxed luma gate ($L_{\max} = 150$).

\textbf{BRU} (\emph{Brukunga, South Australia}): Darker soil and mixed
understorey vegetation. Higher background visual complexity; canonical shadow
luma gate applies ($L_{\max} = 71$).

Tiles of $400 \times 400$\,px are extracted from full orthomosaics. After
augmentation (horizontal/vertical flips, $90°$ rotations), the dataset
comprises 74 training tiles (1,192 annotations) and 20 validation tiles
(301 annotations) across both sites. Domain breakdown is shown in
Table~\ref{tab:dataset}.

\begin{table}[h]
  \centering
  \caption{Dataset statistics after augmentation.}
  \label{tab:dataset}
  \small
  \begin{tabular}{lrrrr}
    \toprule
    & \multicolumn{2}{c}{Train} & \multicolumn{2}{c}{Val} \\
    \cmidrule(lr){2-3}\cmidrule(lr){4-5}
    Site & Tiles & Annots. & Tiles & Annots. \\
    \midrule
    WON & 63 & 1,075 & -- & 258 \\
    BRU &  11 &   117 & -- &  43 \\
    \textbf{Total} & \textbf{74} & \textbf{1,192} & \textbf{20} & \textbf{301} \\
    \bottomrule
  \end{tabular}
\end{table}

\subsection{Evaluation Metric}

We report mean Average Precision at IoU threshold 0.4 (mAP@0.4). The looser
threshold (vs.\ the standard 0.5) is appropriate for aerial TCD where
ground-truth boxes are loosely drawn by human annotators: a prediction that
tightly circumscribes the canopy may fall below IoU=0.5 against an
annotation that includes adjacent shadow. Per-domain mAP (mAP$_{\text{WON}}$,
mAP$_{\text{BRU}}$) is reported separately to surface domain-specific effects.

\subsection{Training Details}

All models fine-tune from a DeepForest NEON checkpoint using the Adam
optimiser with $\text{lr}=0.001$, batch size 16, and early stopping with
patience 10 on validation mAP@0.4 (up to 50 epochs). Training runs on a
single NVIDIA H100 GPU via Modal cloud compute. Data augmentation is disabled
when the shadow channel is active to preserve shadow map consistency with the
RGB image.

\subsection{Ablation Conditions}

We design four training conditions to isolate the contribution of each
component:

\begin{description}
  \item[\textbf{A --- Baseline}] RGB only, no shadow information.
  \item[\textbf{D --- Directional shadow channel}] Shadow map $\mathbf{M}$
        computed with the true sun azimuth $\phi$ appended as 4th channel.
  \item[\textbf{E --- Isotropic shadow channel}] Shadow map computed by
        averaging the directional gradient over 8 equally-spaced azimuths
        (no direction vector used), appended as 4th channel. Identical
        pipeline to D; only the direction is removed. This is the critical
        ablation separating the contribution of the shadow \emph{map} from
        the shadow \emph{direction}.
  \item[\textbf{F --- Shadow input only}] RGB discarded; shadow map tiled
        $3\times$ across channels. Tests whether shadow alone encodes
        sufficient tree-location signal for detection.
        \textit{[Pending]}
\end{description}

% ── 5. Results ────────────────────────────────────────────────────────────────
\section{Results}

\subsection{Main Ablation}

Table~\ref{tab:ablation} reports validation mAP@0.4 across all conditions.
Run D (directional shadow channel) achieves \textbf{0.648} overall mAP vs.\
0.620 for the RGB-only baseline, a gain of +4.5\%. Run E (isotropic shadow
channel, no direction) achieves \textbf{0.650}, statistically indistinguishable
from D. Both shadow conditions converge to approximately the same final
performance across all domains.

\begin{table}[h]
  \centering
  \caption{Ablation results on the validation set (mAP@0.4).
           $\dagger$: pending.}
  \label{tab:ablation}
  \small
  \begin{tabular}{llccc}
    \toprule
    Run & Input & mAP & mAP$_{\text{WON}}$ & mAP$_{\text{BRU}}$ \\
    \midrule
    A & RGB only          & 0.620 & 0.637 & 0.539 \\
    D & RGB + dir.\ shadow  & 0.648 & 0.675 & 0.560 \\
    E & RGB + iso.\ shadow  & \textbf{0.650} & \textbf{0.677} & 0.544 \\
    F & Shadow only        & $\dagger$ & $\dagger$ & $\dagger$ \\
    \bottomrule
  \end{tabular}
\end{table}

The improvement is larger on WON (+6.3\% for D, +6.3\% for E) than BRU
(+3.9\% for D, +0.9\% for E), consistent with WON's sparser canopy and
more prominent shadow contrast on pale substrate. BRU performance is more
variable across runs, likely due to the smaller number of BRU validation
tiles (43 annotations across 5 tiles).

\subsection{Qualitative Analysis}

Figure~\ref{fig:qual} shows representative predictions from runs A and D
on challenging WON tiles. The shadow channel primarily recovers missed trees
(false negatives) at the boundary of shadow-occluded regions, where the RGB
signal alone is ambiguous. False positive rates are broadly stable, suggesting
the shadow channel acts as a recall booster rather than a precision-precision
trade-off.

\begin{figure}[t]
  \centering
  % TODO: replace with actual qualitative figure from review_predictions.py
  \fbox{\rule{0pt}{5cm}\rule{\linewidth}{0pt}}
  \caption{Predictions on WON validation tiles. Left: RGB baseline (A) ---
           green: GT, blue: TP, red: FP, orange: FN (shadow panel).
           Right: shadow-conditioned model (D). Shadow channel recovers
           trees obscured in RGB alone (orange $\to$ blue).}
  \label{fig:qual}
\end{figure}

\subsection{Direction vs.\ Darkness}

The equivalence of runs D and E (directional vs.\ isotropic shadow map) is the
key negative result. Despite the isotropic map being visually noisier ---
activating on any locally-dark region regardless of cause --- it provides the
same detection improvement as the precisely directional map. We interpret this
as follows: the shadow channel improves performance because it explicitly flags
\emph{where} occlusion occurs, allowing the backbone to adapt its feature
statistics for shadowed crowns. The \emph{direction} of the shadow (which
encodes geometric information about relative crown-to-shadow displacement) is
not exploited by the Faster R-CNN architecture in this configuration.

This is an important distinction for the paper's contribution framing: the
novelty of the shadow direction prediction model (Stage 1) lies in enabling
shadow map generation without flight metadata, not in providing architectural
signal that is presently exploited by the detector.

\subsection{Shadow Map Quality}

Figure~\ref{fig:shadow_maps} illustrates the qualitative difference between
directional and isotropic shadow maps on representative WON and BRU tiles.
The directional map produces compact, elongated blobs localised to tree cast
shadows; the isotropic map activates more diffusely on all locally-dark
regions including canopy interior and dark soil patches.

\begin{figure}[t]
  \centering
  % TODO: replace with actual figure from review_shadow_params.py
  \fbox{\rule{0pt}{5cm}\rule{\linewidth}{0pt}}
  \caption{Comparison of shadow maps. Each row: RGB (left), directional map
           (centre), isotropic map (right) at matched luma$_{\max}$ per
           domain. Directional maps are sparser and better localised to
           cast shadow regions. Colour scale: inferno, [0, 1].}
  \label{fig:shadow_maps}
\end{figure}

% ── 6. Discussion ─────────────────────────────────────────────────────────────
\section{Discussion}

\paragraph{Why does the direction not matter?}
The Faster R-CNN backbone processes the shadow map as a spatial intensity
channel. The directional information is encoded implicitly in the elongated
shape of shadow blobs in the map, but the backbone has no mechanism to
interpret that shape as a \emph{directional offset}. The network learns
``dark blob $\approx$ tree present nearby'' without recovering the explicit
geometric relationship $\text{crown} = \text{shadow} - d \cdot \hat{\mathbf{s}}$.

\paragraph{Towards direction-aware architectures.}
Two architectural directions could unlock the shadow direction vector's
geometric value:
\begin{enumerate}
  \item \textbf{Shadow-offset RoI sampling}: augment RoIAlign to sample
        features at positions offset by $d\hat{\mathbf{s}}$ for a range of
        shadow distances $d$, providing the detection head with crown-side
        evidence from the shadow direction. The offset is initialised from
        $\hat{\mathbf{s}}$ rather than learned from scratch.
  \item \textbf{Learned spatial warp}: a lightweight spatial transformer
        network~\cite{jaderberg2015spatial} takes $\hat{\mathbf{s}}$ and
        $\mathbf{M}$ as input and predicts a flow field that shifts shadow
        evidence toward likely crown positions before the backbone processes
        the image. Zero-initialised to identity at training start.
\end{enumerate}
Both approaches are deferred to future work; the ablation infrastructure
(directional vs.\ isotropic conditions F, E) is already in place for
a clean experimental test.

\paragraph{Limitations.}
The validation set is small (20 tiles, 301 annotations) and
domain-imbalanced (WON-heavy). The +4.8\% improvement should be confirmed on
a larger held-out test set before operational deployment. The shadow-only
condition (run F) is pending; its result will determine whether shadow is a
complementary or near-sufficient detection cue in this domain.

% ── 7. Conclusion ─────────────────────────────────────────────────────────────
\section{Conclusion}

We presented a three-stage pipeline for tree crown detection and carbon
estimation in Australian arid rangeland from 0.1\,m\,px$^{-1}$ UAV imagery.
The central contribution is the exploitation of cast shadows as a detection
cue in a domain where RGB appearance alone is insufficient: pale sandy
substrate and sparse canopy create strong visual ambiguity that existing
SOTA detectors fail to resolve.

Stage 1 predicts sun azimuth directly from imagery via a fine-tuned ResNet34,
enabling shadow map generation without flight metadata. Stage 2's shadow-conditioned
Faster R-CNN achieves +4.8\% mAP@0.4 over the RGB-only baseline by appending
the shadow probability map as a fourth input channel. Stage 3 provides per-crown
segmentation via SAM for downstream biomass and carbon estimation.

The ablation study reveals that detection gains are attributable to the shadow
map's role as an \emph{occlusion-awareness signal} rather than to the
directional information encoded in the shadow vector. This is a clean negative
result that motivates future architectures capable of exploiting shadow
geometry explicitly.

% ── References ────────────────────────────────────────────────────────────────
\bibliographystyle{plain}
\bibliography{refs}

\end{document}
